% docs/manual/latex/front/user_reading.tex - front matter of the user manual (PDF only): how to read this book.
% Hand-written LaTeX (no Markdown counterpart); \input by the main .tex of the book.
%   Compile: cmake --build <build> --target pdf, or python3 -I tools/build_docs_pdf.py --source-root . --build-dir DIR
%   (figures, then xelatex passes + makeindex -s docs/common/latex/simcpp.ist); by hand see docs/README.md.
%   XeLaTeX; class docs/common/latex/simcpp.cls (TeX Live packages listed there); fonts TeX Gyre Pagella,
%   TeX Gyre Heros, DejaVu Sans Mono. The figures are built first from docs/{manual,translation}/figures/*.tex
%   (preview.sty) into PDFs.
\section*{How to read this book}
The manual is in four parts. \textbf{Getting started} (chapters~\ref{user_manual-01u5fintroduction}
to~\ref{user_manual-03u5fquicku5fstart}) is enough to install SIMC++ and make a first run. \textbf{Configuring a
simulation} explains the input files and contains the complete parameter reference
(section~\ref{user_manual-04u5finputs:46-parameter-reference}), then what each reaction and each spectrometer does.
\textbf{Results and reproducibility} describes every output file, the random-number generators, the
legacy-compatible mode and multi-threaded runs. \textbf{Help} collects the common messages and problems. The appendix
is a transcript of every command of the manual, run on the version printed on the title page.
